\scanpage{38}
A Moyennant l'isom $(\alpha)$, l'opération $L_Y$ s'identifie à $L_X$ sur le morceau
\[
\coprod_{\substack{i+j\leq n-2\\i,j\geq0}}L_X^jP^j(X),
\]
et son noyau est le morceau
\[
\coprod_{\substack{\deleted{i+j}\\i+j=n-1\\i,j\geq0}}\uline{L_X^jP^j(X)}\oplus E^{(n-1)}
=\coprod_{i\geq n+1}P^j(X)\oplus E^{(n-1)}
\]
(qui n'est d'ailleurs autre que $\displaystyle\coprod_{\substack{\deleted{\unsure{i\geq}}\\n\leq i\leq2(n-1)}}P^i(Y)$, comme il se doit). L'opération $\Lambda_Y$ est nulle sur $E^{n-1}$, et coïncide sur le morceau
\[\coprod_{\substack{i+j=n-1\\i,j\geq0}}L_X^jP^i(X)\]
avec $\Lambda_X$. On retrouve \deleted{l'injectivité}\inserted{que le noyau} de $\Lambda_Y$ \deleted{sur} est
\[
\coprod_{0\leq i\leq n-1}P^i(X)+E^{n-1}\simeq\coprod_{0\leq i\leq n+1}P^i(Y),
\]
\deleted{\unsure{donc}} et que $\Lambda_Y$ est injectif sur le reste
\[
=\coprod_{\substack{i+j\leq n-2\\\deleted{i\geq0,j>0}\\i,j\geq0}}L_X^{1+j}P^i(X).
\]
On en déduit
\figure{\begin{tikzpicture}[x=1cm,y=1cm]
\node[anchor=north west,inner sep=0pt,text width=14cm] at (0,0) {\begin{minipage}{14cm}
\deleted{
\[
\left\{\begin{aligned}
\deleted{\unsure{\Lambda_Y\varphi^*=\Lambda_X}}\varphi^*\Lambda_X&=\Lambda_Y\varphi^*+\\
\varphi_*\Lambda_Y&=\Lambda_X\varphi_*
\end{aligned}\right.
\]
Pour la \deleted{dernière}\inserted{prem.} relation, on la vérifie d'abord sur les $P^j(X)$ \inserted{$=$} pour $j\geq n$, \inserted{i.e. les $L_X^{n-i}P^i(X)$ pour $i\leq n-1$)} \deleted{le premier est} y est nul car $\varphi^*$ l'est, le deuxième
}
\end{minipage}};
\draw (-.1,.08)--(8.2,.35);
\draw (-.1,.08)..controls(.05,-.6)and(.1,-1.3)..(-.7,-1.35)--(-.7,-4.35)--(7.4,-4.1)..controls(8.1,-4.1)and(8.1,-3.95)..(8.1,-3.2)--(13.2,-3.05);
\draw (8.5,.5)..controls(9.1,.7)and(9.1,-.7)..(10,-1)..controls(11.1,-1.3)and(13.2,-.9)..(14,-1.2);
\draw (2.7,-1)--(4.8,-3.9);
\draw (5.2,0)--(7.6,-3.1);
\draw (8.8,-.65)--(10.7,-2.85);
\end{tikzpicture}}

\deleted{$\varphi^*\Lambda_X=$}
\[
\left\{\begin{aligned}
\Lambda_Y\varphi^*&=\varphi^*\Lambda_X\left(\operatorname{id}_X-\sum_{n+1\leq j\leq2n}p_X^j\right)
=\left(\operatorname{id}_Y-\sum_{\substack{\deleted{\unsure{n+1}}\\n-1\leq j\leq2(n-1)}}p_X^j\right)\varphi^*\Lambda_X\\
\varphi^*\Lambda_Xp_X^j&=p_Y^{j-2}\varphi^*\Lambda_X\deleted{=p_Y^{j-2}\Lambda_Y\varphi^*}
\qquad\text{si }j\geq n+1.
\end{aligned}\right.
\]
\margin{\[\Lambda_Y^m\varphi^*=\varphi^*\Lambda_X^m(\operatorname{id}-\bar\Sigma)\quad\Longleftarrow\]
pour tt $m$}
\[
\left\{\begin{aligned}
\varphi_*\Lambda_Y&=\left(\operatorname{id}_X-\sum_{0\leq i\leq n-1}p_X^i\right)\Lambda_X\varphi_*
=\Lambda_X\varphi_*\left(\operatorname{id}_Y-\sum_{0\leq i\leq n-1}p_Y^i\right)\\
p_X^i\Lambda_X\varphi^*&=\Lambda_X\varphi^*p_Y^i=p_X^i\varphi_*\Lambda_Y\qquad\text{si }0\leq i\leq n-1.
\end{aligned}\right.
\]
\mspage{37}
